\section{CÁC SỐ ĐẶC TRƯNG ĐO ĐỘ PHÂN TÁN}
\subsection{TÓM TẮT LÝ THUYẾT}
\subsubsection{Khoảng biến thiên và khoảng tứ phân vị}
Sắp xếp mẫu số liệu theo thứ tự không giảm. Ta có
\begin{itemize}
	\item [\iconMT] \indam{Khoảng biến thiên:} Kí hiệu là $R$, là hiệu số giữa giá trị lớn nhất và giá trị nhỏ nhất trong
	mẫu số liệu.
	\item [\iconMT] \indam{Khoảng tứ phân vị:} Kí hiệu là $ \triangle_Q $, là hiệu số giữa tứ phân vị thứ ba và tứ phân vị thứ nhất, tức là
	\boxmini{$ \triangle_ Q=Q_3-Q_1 $}
	\begin{tcolorbox}[colframe=orange,colback=white,boxrule=0.2mm]
		\indamm{Chú ý:}
		\begin{itemize}
			\item  Khoảng biến thiên dùng để đo độ phân tán của mẫu số liệu. Khoảng biến thiên càng lớn thì mẫu số liệu càng phân tán.
			\item Khoảng tứ phân vị cũng là một số đo độ phân tán của mẫu số liệu. Khoảng tứ phân vị càng lớn thì mẫu số liệu cảng phân tán.
			\item Về bản chất, khoảng tứ phân vị là khoảng biến thiên của $50 \%$ số liệu chính giữa của mẫu số liệu đã sắp xếp.
			\item Một số tài liệu gọi khoảng biến thiên là biên độ và khoảng tứ phân vị là độ trải giữa.
		\end{itemize}
	\end{tcolorbox}
\end{itemize}
\subsubsection{Phương sai và độ lệch chuẩn}
\begin{itemize}
	\item [\iconMT] \indam{Phương sai:} Để đo mức độ phân tán (so với số trung bình cộng) ta dùng phương sai $s^{2}$. Cách tính như sau:
	\begin{itemize}
		\item
		Với mẫu số liệu kích thước $\mathrm{n}$ là $\left\{x_{1}, x_{2}, \ldots, x_{n}\right\}$ thì
		\boxmini{$s^{2}=\dfrac{\left(x_{1}-\overline{x}\right)^{2}+\left(x_{2}-\overline{x}\right)^{2}+\ldots+\left(x_{n}-\overline{x}\right)^{2}}{n}=\dfrac{1}{n}\left(x_1^2+x_2^2+\cdots + x_n^2\right)-\overline{x}^2$}
		\item
		Với mẫu số liệu được cho bởi bảng phân bố tần số thì
		\boxmini{$s^{2}=\dfrac{n_1\left(x_{1}-\overline{x}\right)^{2}+n_2\left(x_{2}-\overline{x}\right)^{2}+\ldots+n_k\left(x_{k}-\overline{x}\right)^{2}}{n}=\dfrac{1}{n}\left(n_1x_1^2+n_2x_2^2+\cdots + n_kx_k^2\right)-\overline{x}^2$}
		Trong đó $n_k$ là tần số của giá trị $x_k$ và $n_1+n_2+ \cdots + n_k=n$.
	\end{itemize}
	\item [\iconMT] \indam{Độ lệch chuẩn:} Căn bậc hai của phương sai gọi là độ lệch chuẩn, kí hiệu là $s$. Ta có $s=\sqrt{s^2}$.
	\begin{tcolorbox}[colframe=orange,colback=white,boxrule=0.2mm]
		\indamm{Chú ý:}
		\begin{itemize}
			\item  Phương sai và độ lệch chuẩn càng lớn thì độ phân tán của các số liệu thống kê càng lớn.
			\item  Phương sai $s^2$ và độ lệch chuẩn $s$ đều được dùng để đánh giá mức độ phân tán của các số liệu thống kê (so với số trung bình cộng). Nhưng khi cần chú ý đến đơn vị đo thì ta dùng $s$ vì $s$ có cùng đơn vị đo với dấu hiệu được nghiên cứu.
		\end{itemize}
	\end{tcolorbox}
\end{itemize}
\subsubsection{Phát hiện số liệu bất thường hoặc không chính xác bằng biểu đồ hộp}
Trong mẫu số liệu thống kê, có khi gặp những giá trị quá lớn 
hoặc quá nhỏ so với đa số các giá trị khác. Những giá trị này được gọi 
là \indamm{giá trị bất thường}. Ta có thể dùng biểu đồ hộp để phát hiện các giá trị bất thường 
này.
\begin{center}
	\begin{tikzpicture}[>=stealth, font=\footnotesize, line join=round, line cap=round,transform shape,scale=1]
		\begin{scope}[scale=1]
			%% Vẽ hộp
			\foreach \t in {7, 8.5, 2, 11} \draw (\t,0.5)--(\t,2);
			\draw (5.2,0.5)--(8.5,0.5) (5.2,2)--(8.5,2);
			\draw (2,1.25)--(11,1.25);
			\fill[magenta!20!white] (5.2,0.5) rectangle (8.5,2);
			\draw[color=yellow, thick] (7,0.5)--(7,2);
			%% ghi các giá trị phía dưới
			\foreach \e/\f  in {2/\big(Q_1-1{,}5 \cdot  
				\Delta_{Q}\big),7/Q_2,5.2/Q_1,8.5/Q_3,11/\big(Q_3+1{,}5 \cdot 
				\Delta_{Q}\big)}
			{
				\node[shift={(-90:4mm)}] at (\e,0.5){$\f$};
			}
			%% Vẽ giá rị bất thường
			\foreach \btd  in {0.2,0.6}	
			\draw[fill=orange] (2-\btd,1.25) circle (0.08);
			\foreach \btt  in {0.2,0.4}	
			\draw[fill=orange]  (11+\btt,1.25) circle (0.08);
			\draw [thick,->](2-0.2,2.2)--(2-0.2,1.35);
			\draw [thick,->](11+0.4,2.2)--(11+0.4,1.35);
			\node [text = red,
			fill = cyan!30!white,rounded corners=8pt] at (2 -0.5,2.5){Các giá trị bất thường};
			\node [text = red,
			fill = cyan!30!white,rounded corners=8pt]  at (11 +0.5,2.5){Các giá trị bất thường};
			%%% Vẽ delta Q
			\draw (5.2,2 +0.2)--node[above]{$\Delta_{Q}$}(8.5,2+0.2);
			\draw (5.2,2 +0.1)--(5.2,2 +0.3) (8.5,2 +0.1)--(8.5,2 +0.3);
		\end{scope}
		\draw (-1,-1) rectangle (14,2+1);
	\end{tikzpicture}
\end{center}
Các số liệu \indamm{lớn hơn} $Q_3+1{,}5 \cdot 
\Delta_{Q}$ hoặc \indamm{bé hơn} $Q_1-1{,}5 \cdot 
\Delta_{Q}$ được xem là giá trị bất thường.
\subsection{RÈN LUYỆN KĨ NĂNG GIẢI TOÁN}
\begin{dang}{Tìm khoảng biến thiên của mẫu số liệu}
	\begin{itemize}
		\item [\iconMT] \indam{Phương pháp:}  Tìm giá trị lớn nhất $M$ và giá trị nhỏ nhất $m$ của mẫu số liệu. Khi đó, koảng biến thiên của mẫu số liệu là $R=M-m$.
	\end{itemize}
\end{dang}

\begin{vd}
	Tuổi thọ trung bình người dân của $ 11 $ nước được thống kê như sau:
	\begin{longtable}{p{0,7cm}p{0.7cm}p{0.7cm}p{0.7cm}p{0.7cm}p{0.7cm}p{0.7cm}p{0.7cm}p{0.7cm}p{0.7cm}p{0.7cm}}
		69 & 77 & 75 & 83 & 65 & 75 & 74 & 68 & 73 & 72 & 71
	\end{longtable}
	Hãy tìm khoảng biến thiên của mẫu số liệu trên.
	\loigiai{
		Khoảng biến thiên của mẫu số liệu là $R=83-65=18$.		
	}
\end{vd}\dongcham{2}

\begin{vd}
	Điểm thi học kì 2 môn Toán và Ngữ văn của một nhóm học sinh được ghi lại như sau
	\begin{longtable}{|l|c|c|c|c|c|c|c|c|c|}
		\hline 
		Toán & 9 & 8{,}5 & 7 & 6{,}3 & 5 & 9{,}5 & 8 \\ 
		\hline 
		Ngữ văn & 6 & 6{,}5 & 8 & 7{,}3 & 5{,}5 & 8{,}3 & 6{,}5 \\ 
		\hline 
	\end{longtable}
	Hãy tìm khoảng biến thiên của hai mẫu số liệu trên. Từ đó chỉ ra mẫu số liệu có độ phân tán lớn hơn.
	\loigiai{
		Khoảng biến thiên của mẫu số liệu ``Toán'' là $ R_1=9{,}5-5=4{,}5 $.\\
		Khoảng biến thiên của mẫu số liệu ``Ngữ văn'' là $ R_2=8{,}3-5{,}5=2{,}8 $.\\
		Do $R_1>R_2$ nên mẫu số liệu ``Toán'' có độ phân tán lớn hơn mẫu số liệu ``Ngữ văn''.
	}
\end{vd}
\begin{dang}{Tìm khoảng tứ phân vị của mẫu số liệu}
	\begin{itemize}
		\item [\iconMT] \indam{Phương pháp:}  Ta xác định các tứ phân vị $Q_1$, $Q_2$, $Q_3$. Khi đó khoảng tứ phân vị là $$\triangle_Q=Q_3-Q_1.$$
	\end{itemize}
\end{dang}
\begin{vd}
	Mẫu số liệu sau đây cho biết cân nặng của 10 trẻ sơ sinh (đơn vị kg)
	\begin{itemize}
		\item [] 2,977 \quad 3,155 \quad 3,920 \quad 3,412 \quad 4,236
		\item [] 2,593 \quad 3,270 \quad 3,813 \quad 4,042 \quad 3,387
	\end{itemize}
	Hãy tính khoảng biến thiên, khoảng tứ phân vị cho mẫu số liệu này.
	\loigiai{
		Trước hết, ta sẽ sắp xếp mẫu số liệu theo thứ tự không giảm:
		$$2,593 \quad 2,977 \quad 3,155 \quad 3,270 \quad 3,387 \quad 3,412 \quad 3,813 \quad 3,920 \quad 4,042 \quad 4,236$$
		Khi đó
		\begin{itemize}
			\item [$\bullet$] Khoảng biến thiên là $R=4,236-2,593=1,643$ .
			\item [$\bullet$] $Q_2=3,3995$ ; $Q_1=3,155$ ; $Q_3=3,920$ nên khoảng tứ phân vị là $\Delta_Q=Q_3-Q_1=0,765$.
		\end{itemize}
	}
\end{vd}\dongcham{8}

\begin{vd}
	Nhiệt độ trung bình các tháng trong năm 2019 của hai tỉnh Lai Châu và Lâm Đồng được ghi lại trong bảng sau:
	\begin{center}
		\begin{tikzpicture}
			\matrix[matrix of nodes,nodes in empty cells,
			row sep=-\pgflinewidth,column sep=-\pgflinewidth,
			nodes={minimum height=8mm,minimum width=11mm,draw=black,anchor=center},
			column 1/.style={nodes={minimum width=24mm,color=black,fill=cyan!10}},
			row 1/.style={nodes={fill=cyan!10}},
			row 2/.style={nodes={minimum height=12mm}},
			row 3/.style={nodes={minimum height=12mm}},
			]{
				Tháng &1&2&3&4&5&6&7&8&9&10&11&12\\ 
				\node[align=center]{Lai Châu}; &14,8&18,8&20,3&23,5&24,7&24,2&23,6&24,6&22,7&21,0&18,6&14,2\\
				\node[align=center]{Lâm Đồng}; &16,3&17,4&18,7&19,8&20,2&20,3&19,5&19,3&18,6&18,5&17,5&16,0\\
			};
		\end{tikzpicture}
	\end{center}
	\begin{tasks}(1)
		\task Tìm khoảng biến thiên và khoảng tứ phân vị của mẫu số liệu trên.
		\task Hãy cho biết trong một năm, nhiệt độ địa phương nào ít thay đổi hơn?
	\end{tasks}
	\loigiai{
		Xếp hai dãy số liệu trên theo thứ tự không giảm:
		\begin{itemize}
			\item [] Số liệu của tỉnh Lai Châu: 14,2 \quad 14,8 \quad 18,6 \quad 18,8 \quad 20,3 \quad 21,0 \quad 22,7 \quad 23,5 \quad 23,6 \quad 24,2 \quad 24,6 \quad 24,7.
			\item [] Số liệu của tỉnh Lâm Đồng: 16,0 \quad 16,3 \quad 17,4 \quad 17,5 \quad 18,5 \quad 18,6 \quad 18,7 \quad 19,3 \quad 19,5 \quad 19,8 \quad 20,2 \quad 20,3.
		\end{itemize}
		
		\begin{enumerate}[a)]
			\item Đối với tỉnh Lai Châu thì
			\begin{itemize}
				\item [$\bullet$] $R=24,7-14,2=10,5$;
				\item [$\bullet$] $Q_1=\dfrac{18,6+18,8}{2}=18,7$ và $Q_3=\dfrac{23,6+24,2}{2}=23,9$. Suy ra $\Delta_Q=Q_3-Q_1=5,2$.
			\end{itemize}
			Đối với tỉnh Lâm Đồng thì
			\begin{itemize}
				\item [$\bullet$] $R=20,3-16,0=4,3$;
				\item [$\bullet$] $Q_1=\dfrac{17,4+17,5}{2}=17,45$ và $Q_3=\dfrac{19,5+19,8}{2}=19,65$. Suy ra $\Delta_Q=Q_3-Q_1=2,2$.
			\end{itemize}
			\item So sánh khoảng biến thiên hoặc khoảng tứ phân vị, ta thấy nhiệt độ tỉnh Lâm Đồng ít thay đổi hơn (độ phân tán thấp hơn Lai Châu).
		\end{enumerate}
	}
\end{vd}\dongcham{8}
\begin{dang}{Tìm phương sai và độ lệch chuẩn của mẫu số liệu}
	\begin{itemize}
		\item [\iconMT] \indam{Tìm phương sai $s^2$:}
		\begin{itemize}
			\item [$\bullet$] Tính số trung bình  $\overline{x}$ của mẫu số liệu.
			\item [$\bullet$] Thay vào công thức: $s^{2}=\dfrac{\left(x_{1}-\overline{x}\right)^{2}+\left(x_{2}-\overline{x}\right)^{2}+\ldots+\left(x_{n}-\overline{x}\right)^{2}}{n}$
		\end{itemize}
		\item [\iconMT] \indam{Tìm độ lệch chuẩn $s$:} Tính phương sai $s^2$ của mẫu số liệu, suy ra độ lệch chuẩn $s=\sqrt{s^2}$.
	\end{itemize}
\end{dang}

\begin{vd}
	Bảng dưới đây thống kê tổng số giờ nắng trong năm 2019 theo từng tháng được đo bởi hai trạm quan sát khí tượng đặt ở Tuyên Quang và Cà Mau.
	\begin{center}
		\begin{tikzpicture}
			\matrix[matrix of nodes,nodes in empty cells,
			row sep=-\pgflinewidth,column sep=-\pgflinewidth,
			nodes={minimum height=8mm,minimum width=11mm,draw=black,anchor=center},
			column 1/.style={nodes={minimum width=28mm,color=black,fill=cyan!10}},
			row 1/.style={nodes={fill=cyan!10}},
			row 2/.style={nodes={minimum height=12mm}},
			row 3/.style={nodes={minimum height=12mm}},
			]{
				Tháng &1&2&3&4&5&6&7&8&9&10&11&12\\ 
				\node[align=center]{Tuyên Quang}; &25&89&72&117&106&117&156&203&227&146&117&145\\
				\node[align=center]{Cà Mau}; &180&223&257&245&191&111&141&134&130&122&157&173\\
			};
		\end{tikzpicture}
	\end{center}
	\begin{tasks}(1)
		\task Hãy tính phương sai và độ lệch chuẩn của dữ liệu từng tỉnh.
		\task Nêu nhận xét về sự thay đổi tổng số giờ nắng theo từng tháng ở mỗi tỉnh.
	\end{tasks}
	\loigiai{
		\begin{enumerate}[a)]
			\item Tại Tuyên Quang: 
			\begin{itemize}
				\item [$\bullet$] Số giờ nắng trung bình mỗi tháng là $\overline{x}=\dfrac{395}{3}$.
				\item [$\bullet$] Phương sai:
				$$s^2=\dfrac{1}{12}(x_1^2+x_2^2+\cdots + x_{12}^2)-\overline{x}^2\approx 2921,22.$$
				\item [$\bullet$] Độ lệch chuẩn $s=\sqrt{s^2}\approx 54,05$.
			\end{itemize}
			Tại Cà Mau:
			\begin{itemize}
				\item [$\bullet$] Số giờ nắng trung bình mỗi tháng là $\overline{x}=172$.
				\item [$\bullet$] Phương sai:
				$$s^2=\dfrac{1}{12}(x_1^2+x_2^2+\cdots + x_{12}^2)-\overline{x}^2\approx 2183.$$
				\item [$\bullet$] Độ lệch chuẩn $s=\sqrt{s^2}\approx 46,72$.
			\end{itemize}
			\item Theo kết quả của phương sai và độ lệch chuẩn thì tổng số giờ nắng ở Tuyên Quang theo từng tháng có sự biến động lớn hơn ở Cà Mau.
		\end{enumerate}
	}
\end{vd}\dongcham{7}
\begin{vd}%[0D5B4]
	Sản lượng lúa (đơn vị là tạ) của $40$ thửa ruộng thí nghiệm có cùng diện tích được trình bày trong bảng tần số dưới đây:
	\begin{center}
		\begin{tabular}{|c|c|c|c|c|c|c|}
			\hline 
			Sản lượng ($x$) & $20$ & $21$ & $22$ & $23$ & $24$ & Tổng \\ 
			\hline 
			Tần số ($n$) & $5$ & $8$ & $11$ & $10$ & $6$ & $n=40$ \\ 
			\hline 
		\end{tabular} 
	\end{center}
	\begin{tasks}(1)
		\task Tính sản lượng trung bình của $40$ thửa ruộng?
		\task Tính phương sai và độ lệch chuẩn.
	\end{tasks}
	\loigiai{
		\begin{enumerate}[a)]
			\item Số trung bình của sản lượng của $40$ thửa ruộng là:
			$$\overline{x}=\dfrac{5\cdot 20+8\cdot 21 +11\cdot 22+10\cdot 23+6\cdot 24}{40}=22,1 \text{ tạ }.$$
			\item Tính phương sai:
			\begin{eqnarray*}
				s^2
				&=&\dfrac{1}{n} \sum\limits_{i=1}^{5} n_i\left(x_i-\overline{x}\right)^2\\
				&=&\dfrac{1}{40}\left[ 5(20-22,1)^2+8(21-22.1)^2+11(22-22,1)^2+10(23-22,1)^2+6(24-22,1)^2\right]\\
				&=& \dfrac{6160}{4000}=1,54.
			\end{eqnarray*}
			Tính độ lệch chuẩn: $s=\sqrt{s^2}=\sqrt{1,54}\approx 1,24.$
		\end{enumerate}
	}
\end{vd}\dongcham{7}

\begin{dang}{Tìm các số liệu bất thường của mẫu số liệu}
	Để tìm các số liệu bất thường của một mẫu số liệu ta thực hiện các bước giải như sau:
	\begin{itemize}
		\item Tìm khoảng tứ phân vị $\Delta_{Q}=Q_3 - Q_1$ của mẫu số liệu.
		\item Xác định đoạn $\left[Q_1-1{,}5 \cdot 
		\Delta_{Q}; Q_3+1{,}5 \cdot 
		\Delta_{Q} \right]$.
		\item Tìm các số liệu \textbf{không} thuộc đoạn $\left[Q_1-1{,}5 \cdot	\Delta_{Q}; Q_3+1{,}5 \cdot \Delta_{Q} \right]$ là \indamm{các số liệu bất thường}
	\end{itemize}
\end{dang}

\begin{vd}%[0D5B3-5]
	Một cảnh sát giao thông ghi tốc độ (đơn vị: km/h) của $25$ chiếc xe qua trạm như sau:
	\begin{center}
		\begin{tabular}{ccccccccccccccc}
			20 & 41 & 41 & 80 & 40 & 52 & 52 & 52 & 60 & 55 & 60 & 60 & 62  \\  
			60 & 65 & 60 & 65 & 135 & 70 & 70 & 65 & 75 & 75 & 70 & 55 &  \\  
		\end{tabular} 
	\end{center}
	Hãy tìm các số liệu bất thường trong mẫu số liệu trên.
	\loigiai{
		Sắp xếp các số liệu trong mẫu theo thứ tự không giảm ta có
		\begin{center}
			\begin{tabular}{|c|c|c|c|c|c|c|c|c|c|c|c|c|}
				\hline 
				Tốc độ    & $20$ & $40$& $41$& $52$& $55$& $60$& $62$& $65$& $70$ & $75$& $80$& $135$\\
				\hline 
				Số lần xuất hiện  & $1$ & $1$ & $2 $ & $3$& $2$& $5$& $1$& $3$& $3$& $2$ & $1$& $1$\\
				\hline
			\end{tabular}
		\end{center} 
		Từ bảng số liệu ta tìm được
		\begin{itemize}
			\item [$\bullet$] Số trung vị $Q_2=60$, tứ phân vị dưới $Q_1= 52$, tứ phân vị trên $Q_3= 70$.
			\item [$\bullet$] Khoảng tứ phân vị $\Delta_{Q} = Q_3-Q_1=70 - 52 = 18$.
		\end{itemize}
		Ta có $\left[Q_1-1{,}5 \cdot  \Delta_{Q}; Q_3+1{,}5 \cdot  \Delta_{Q} \right] =\left[25;97\right]$.Số liệu $135$ không thuộc khoảng trên nên nó là số liệu bất thường trong mẫu số liệu.		
	}
\end{vd}\dongcham{10}

\begin{vd}
	Bình dùng đồng hồ đo thời gian để một vật rơi tự do (đơn vị: giây) từ vị trí A đến vị trí B trong 10 lần, cho kết quả như sau:\\
	\begin{tikzpicture}
		\matrix[matrix of nodes,nodes in empty cells,
		row sep=-\pgflinewidth,column sep=-\pgflinewidth,
		nodes={minimum height=8mm,minimum width=14mm,draw=black,anchor=center},
		column 1/.style={nodes={minimum width=25mm,color=black,fill=cyan!10}},
		row 1/.style={nodes={fill=cyan!10}},
		row 2/.style={nodes={minimum height=12mm}},
		]{
			Lần đo &1&2&3&4&5&6&7&8&9&10\\ 
			\node[align=center]{Số giây}; &0,398&0,399&0,408&0,410&0,406&0,405&0,402&0,401&0,290&0,402\\
		};
	\end{tikzpicture}
	Bình nghĩ là giá trị 0,290 ở lần đo thứ 9 không chính xác. Hãy kiểm tra nghi ngờ của Bình.
	\loigiai{Xếp các số liệu trên theo thứ tự không giảm:
		\begin{center}
			0,290 \quad 0,398 \quad 0,399 \quad 0,401 \quad 0,402 \quad 0,402 \quad 0,405 \quad 0,406 \quad 0,408 \quad 0,410
		\end{center}
		\begin{itemize}
			\item [$\bullet$] Xét 5 số liệu đầu ở dãy trên, ta tìm được $Q_1=0,399$.
			\item [$\bullet$] Xét 5 số liệu sau ở dãy trên, ta tìm được $Q_3=0,406$.
		\end{itemize}
		Khoảng tứ phân vị của mẫu số liệu trên là $\Delta_{Q}=Q_3-Q_1=0,007$.\\
		Ta có $\left[Q_1-1{,}5 \cdot  \Delta_{Q}; Q_3+1{,}5 \cdot  \Delta_{Q} \right] =\left[25;97\right]=\big[0,3885;0,4165\big]$. Số liệu 0,290 không thuộc khoảng này, nên nó là số liệu bất thường nên nó không chính xác.
	}
\end{vd}\dongcham{12}
